\documentclass[10pt,a4paper]{amsart}
\usepackage[margin=19mm]{geometry}
\usepackage{amsmath,amssymb,mathtools}
\usepackage[scr=rsfs,cal=euler]{mathalfa}
\usepackage{xfrac}
\usepackage[unicode,hidelinks,bookmarks=false]{hyperref}
\newcommand{\ie}{i.e.\ }
\newcommand{\cf}{cf.\ }
\newcommand{\resp}{respectively\ }
\newcommand{\Subsection}{\subsection}
\providecommand{\nobreakdash}{\nobreak\hbox{-}}
\setlength{\emergencystretch}{2em}
\multlinegap0pt
\usepackage{amssymb}
\usepackage{mathtools}
\usepackage[shortlabels]{enumitem}
\usepackage{tikz}
\newtheorem{proposition}{Proposition}
\title{Perimetric Contractions and Their Iterates in Complete $b$-Metric Spaces}
\author{Mujahid Abbas, Alemayehu G. Negash, Meaza F. Bogale}
\date{Source: arXiv:2606.30954v2 (CC BY 4.0)}
\begin{document}
\maketitle

\noindent\emph{Source and licence.} Perimetric Contractions and Their Iterates in Complete $b$-Metric Spaces. Version arXiv:2606.30954v2 (CC BY 4.0), distributed by arXiv under the Creative Commons Attribution 4.0 International licence, http://creativecommons.org/licenses/by/4.0/, as recorded in the arXiv licence metadata for this exact version and shown on the abstract page https://arxiv.org/abs/2606.30954v2. Full licence notice: \texttt{licenses/arxiv-2606.30954-CC-BY-4.0.txt}.\par\smallskip

\noindent\textit{Editorial scope.} Counted unit: the article's proposition prop:continuity (source0.tex lines 355--358). The article's complete original proof of it is delivered with it (source0.tex lines 360--430, one \texttt{proof} environment of 71 lines). No mathematics is written, repaired or re-derived: the statement and the proof are the article's own lines, in the article's own notation, and no retained line is substituted or re-flowed.  No further source-local context is needed: every cross-reference of every retained line resolves inside the two delivered spans.  Nothing was omitted from any retained span, so the package carries no omission record.  No retained line contains a citation, so the wrapper carries no bibliography and no bibliography entry.  No retained line contains a figure environment, an image-inclusion command or drawing code, so no image is delivered and the package declares no asset dependency.  Editorial formatting only, in addition: an AMS \texttt{amsart} preamble that restates the article's own declarations and macros used by the retained lines, namely the environments \texttt{proposition} on separate counters, together with the standard packages those lines need.  The retained environments print in this document as Proposition 1 in order of appearance, whereas the article prints the same items with the numbering of its own full document.  The complete native source of the article as distributed in the arXiv e-print for this exact version is bundled unchanged as \texttt{sources/204062/source0.tex}.
\begin{proposition}
\label{prop:continuity}
Let $(X,d)$ be a $b$-metric space with parameter $s\geq 1$. If $f:X\to X$ is a contracting perimeters of triangles mapping (CPTM), then $f$ is continuous on $X$.
\end{proposition}

\begin{proof}
Let $x\in X$ be an arbitrary point, and let $(y_{m})_{m\in \mathbb{N}}$ be a sequence in $X$ converging to $x$, so that $\lim_{m\to \infty}d(y_{m},x) = 0$. To prove continuity, it suffices to show that $\lim_{m\to \infty}f(y_{m}) = f(x)$.

Assume for contrary that $f(y_{m})\not\to f(x)$. By definition of non-convergence, there exist an $\epsilon >0$ and a subsequence $(y_{m_k})_{k\in \mathbb{N}}$ such that
\[
d(f(y_{m_k}),f(x))\geq \epsilon \quad \forall k\in \mathbb{N}.
\]
We now divide the proof into following four steps.\newline
\textbf{Step 1}:  We first establish that we can extract a subsequence of $(y_{m_k})$ consisting of pairwise distinct points, all different from $x$.

If $y_{m_k} = x$ for infinitely many indices, then for those indices we would have $d(f(y_{m_k}),f(x)) = d(f(x),f(x)) = 0$, directly contradicting the assumption that $d(f(y_{m_k}),f(x))\geq \epsilon$. Thus, $x$ can appear at most finitely many times in the sequence $(y_{m_k})$. Discarding these finite terms allows us to assume without loss of generality that $y_{m_k}\neq x$ for all $k\in \mathbb{N}$. Next, we claim that any unique value $z\in X\backslash \{x\}$ can appear at most finitely many times in $(y_{m_k})$. Suppose some $z\neq x$ appears infinitely many times. Then there exists a subsequence identically equal to $z$. Because any subsequence of a convergent sequence must share the same limit, this constant subsequence must converge to $x$. By the identity of indiscernibles, this implies $d(z,x) = 0$, yielding $z = x$, which is a contradiction.

Since each unique value occurs only finitely many times across an infinite sequence, the set of values $\{y_{m_k}:k\in \mathbb{N}\}$ must be infinite. Therefore, we can inductively extract a subsequence which we reindex and denote as $(y_{m_i})_{i\in \mathbb{N}}$ such that $y_{m_i}\neq y_{m_j}$ for all $i\neq j$. Thus, the set $\{x\} \cup \{y_{m_i}:i\in \mathbb{N}\}$ consists of mutually distinct points.

\textbf{Step 2}: This step deals with the finite image set case (Pigeonhole Principle).\newline Let $F = \{f(y_{m_i}):i\in \mathbb{N}\}$ be the set of image values under $f$. Suppose on contrary that $F$ is a finite set.

By the Pigeonhole Principle, since the index set $\mathbb{N}$ is infinite and $F$ is finite, there must exist a specific point $c\in X$ and an infinite subsequence of $(y_{m_i})$ (which we still denote as $(y_{m_i})$ for simplicity) such that $f(y_{m_i}) = c$ for all $i\in \mathbb{N}$. As $d(f(y_{m_i}),f(x))\geq \epsilon >0$, $c\neq f(x)$.

Now, select any two distinct indices $i\neq j$. The domain elements $x,y_{m_i},y_{m_j}$ are mutually distinct by construction. Applying the CPTM contractive condition to this triple, we have
\[
P(f(x),f(y_{m_i}),f(y_{m_j}))\leq qP(x,y_{m_i},y_{m_j}).
\]
Substituting $f(y_{m_i}) = f(y_{m_j}) = c$ into the left-hand side, we get
\[
P(f(x),c,c) = d(f(x),c) + d(c,c) + d(c,f(x)) = 2d(f(x),c).
\]
Using the $b$-triangle inequality to expand the domain perimeter on the right-hand side, and noting that $d(y_{m_i},y_{m_j})\leq s[d(y_{m_i},x) + d(x,y_{m_j})]$, the inequality becomes
\[
2d(f(x),c)\leq q(s + 1)\left[d(x,y_{m_i}) + d(x,y_{m_j})\right].
\]
As this holds for all distinct $i$ and $j$, on taking limit as $i,j\to \infty$ and applying the Squeeze Theorem (since $y_{m_i},y_{m_j}\to x$), we obtain that
\[
2d(f(x),c)\leq q(s + 1)[0 + 0] = 0\Rightarrow d(f(x),c) = 0\Rightarrow c = f(x),
\]
a contradiction.

\textbf{Step 3}: This step deals with formal inductive construction for the infinite image case.\newline As $F$ is infinite, we can construct a further subsequence whose images are completely distinct from each other and from $f(x)$. We define this subsequence $(y_{m_{i_k}})_{k \in \mathbb{N}}$ inductively:

\textit{Base Step $(k = 1)$:} Choose $i_{1} = 1$. Since $F$ is infinite, $F\backslash \{f(x)\}$ remains infinite; hence, we can ensure $f(y_{m_{i_1}})\neq f(x)$.\newline \textit{Inductive Step:} Assume we have successfully chosen indices $i_1< i_2< \dots < i_k$ such that the set $\{f(y_{m_{i_1}}),\ldots,f(y_{m_{i_k}})\}$ consists of $k$ distinct values, none of which equal $f(x)$. Consider the tail set $F_{k} = \{f(y_{m_{i}}):i > i_{k}\}$. Removing a finite number of elements from the infinite set $F$ leaves $F_{k}$ infinite. Therefore, $F_{k}$ cannot be contained within the finite set $V_{k} = \{f(x),f(y_{m_{i_1}}),\ldots,f(y_{m_{i_k}})\}$. Thus, there must exist an index $i_{k + 1} > i_{k}$ such that $f(y_{m_{i_{k + 1}}})\notin V_{k}$.

By mathematical induction, we obtain a subsequence where both the domain elements are pairwise distinct and their images are pairwise distinct and omit $f(x)$. To simplify notation, let us re-index this and obtain a subsequence as $(w_{k})_{k \in \mathbb{N}}$.

\textbf{Step 4}: This step is about analytical Squeezing and limit Passing.\newline For any two distinct indices $k \neq l$, the domain triple $(x, w_k, w_l)$ and the image triple $(f(x), f(w_k), f(w_l))$ both consist of mutually distinct points. Applying the perimetric contraction condition yields
\[
P(f(x),f(w_k),f(w_l))\leq qP(x,w_k,w_l).
\]
Expanding the perimeters gives
\[
d(f(x),f(w_k)) + d(f(w_k),f(w_l)) + d(f(w_l),f(x))\leq q[d(x,w_k) + d(w_k,w_l) + d(w_l,x)].
\]
Bounding the right-hand side via $d(w_k, w_l) \leq s[d(w_k, x) + d(x, w_l)]$ leads to
\[
d(f(x),f(w_k)) + d(f(w_k),f(w_l)) + d(f(w_l),f(x))\leq q(s + 1)[d(x,w_k) + d(x,w_l)].
\]
Since all metrics on the left-hand side are strictly non-negative, we can drop the second and third terms to isolate the first term:
\begin{equation}
d(f(x),f(w_k))\leq q(s + 1)[d(x,w_k) + d(x,w_l)].
\label{eq:cont_step4}
\end{equation}
Inequality~\eqref{eq:cont_step4} holds strictly for a fixed index $k$ and any choice of $l \in \mathbb{N} \setminus \{k\}$. To evaluate this behavior as $l \to \infty$, we must account for the fact that the $b$-metric $d(\cdot, \cdot)$ is generally not continuous in both variables simultaneously. To rigorously bypass this structural limitation, we take the limit supremum as $l \to \infty$ on both sides of~\eqref{eq:cont_step4}. 

Since the left-hand side is completely independent of $l$, it behaves as a constant bound relative to the running index. Utilizing the fact that $\lim_{l \to \infty} d(x, w_l) = 0$, Lemma~2.2 justifies passing the limit into the right-hand side, yielding
\[
d(f(x), f(w_k)) \le \limsup_{l \to \infty} q(s + 1) \left[ d(x, w_k) + d(x, w_l) \right] = q(s + 1) d(x, w_k).
\]
Finally, taking the limit as $k \to \infty$ on both sides of this reduced inequality, and noting that $\lim_{k \to \infty} d(x, w_k) = 0$, the standard Squeeze Theorem gives 
\[
\lim_{k \to \infty} d(f(x), f(w_k)) = 0.
\]
This contradicts our  assumption that $d(f(w_k), f(x)) \ge \epsilon > 0$ for all $k \in \mathbb{N}$. Thus, our assumption is wrong, and hence $f$ is continuous at $x$. Since $x \in X$ was arbitrarily selected, $f$ is continuous on $X$.
\end{proof}

\end{document}
